% =============================================================================
% APPENDIX MS: REF-MODELSPACE - Scientific Theory Catalog
% =============================================================================
% Category: REF (Reference)
% Language: English
% Version: 1.0 (January 2026)
% Status: Active
%
% Purpose: Comprehensive catalog of established economic and behavioral theories
%          mapped to their EBF (C, Ψ) restrictions. Enables scientific understanding
%          of how theories relate within the EBF framework.
%
% Target Audience: Scientists, Researchers, Academics
%
% Complements:
%   - DSN (METHOD-DESIGN): For practitioners designing models
%   - DER (FORMAL-LIMITS): For formal proofs of 12 key theories
%   - REG (METHOD-REGISTRY): For storing custom-designed models
%   - SW (METHOD-SWSM): Scientific Writing Structure Model (text structure theories)
%
% =============================================================================

\chapter{REF-MODELSPACE: Scientific Theory Catalog}
\label{app:ref-modelspace}

% =============================================================================
% ABSTRACT
% =============================================================================
\begin{abstract}
This appendix provides a comprehensive catalog of 134 established economic and
behavioral theories mapped to their Evidence-Based Framework (EBF) restrictions.
Each theory is characterized by specific constraints on complementarity ($C_{ij}$),
context ($\Psi$), identity utility ($U^{IDN}$), loss aversion ($\lambda$), and
present bias ($\beta$). The catalog enables scientists to understand how theories
relate within the $(C, \Psi)$ framework, and supports the hybrid workflow where
practitioners design 10C models and then ground them in established literature.
Organized in 13 categories from Classical Economics to Development \& Poverty,
this reference serves as the scientific counterpart to the practitioner-focused
model design workflow (Appendix UNMAPPED_DSN).
\end{abstract}

% =============================================================================
% QUICK REFERENCE
% =============================================================================
\begin{tcolorbox}[colback=yellow!10!white, colframe=orange!75!black,
    title=\textbf{Quick Reference: Theory Catalog}]

\textbf{Total:} 134 models in 13 categories

\textbf{Key Restriction Dimensions:}
\begin{itemize}[nosep]
    \item $C_{ij} = 0$: Separable utilities (Classical)
    \item $C_{ij} \neq 0$: Interdependent utilities (Social Preferences, Games)
    \item $\lambda > 1$: Loss aversion (Prospect Theory, Behavioral Finance)
    \item $\beta < 1$: Present bias (Hyperbolic Discounting)
    \item $U^{IDN} \neq 0$: Identity utility (Identity Economics)
    \item $\Psi$ endogenous: Context evolves (Institutions)
\end{itemize}

\textbf{Usage:} Design model with 10C $\rightarrow$ Find similar theory here $\rightarrow$ Use literature for validation

\end{tcolorbox}

% =============================================================================
% HEADER BLOCK
% =============================================================================
\begin{tcolorbox}[colback=blue!5!white, colframe=blue!75!black,
    title=\textbf{Appendix UNMAPPED_MS: REF-MODELSPACE}]

\begin{tabular}{@{}ll@{}}
\textbf{Appendix:} & MS (Scientific Theory Catalog) \\
\textbf{Category:} & REF (Reference) \\
\textbf{Purpose:} & Map established theories to EBF restrictions \\
\textbf{Audience:} & Scientists, Researchers, Academics \\
\textbf{Version:} & 1.1 \\
\textbf{Models:} & 134 established theories (13 categories) \\
\end{tabular}

\tcblower

\textbf{Core Contribution:} Provides a comprehensive catalog showing how
established economic and behavioral theories are special cases of the
$(C, \Psi)$ framework, enabling scientific understanding and theory comparison.

\textbf{Key Distinction:} EBF does NOT refute these theories—it specifies
their \textit{domain of validity} $\Omega_\mathcal{T}$.

\end{tcolorbox}

% =============================================================================
% APPENDIX SCOPE BOX
% =============================================================================
\begin{tcolorbox}[
    colback=orange!5!white,
    colframe=orange!75!black,
    title={\textbf{Appendix Scope: Ziel / In-Scope / Out-of-Scope / Constraints}},
    fonttitle=\bfseries
]

\textbf{Ziel (Objective):}
\begin{itemize}[nosep]
    \item Catalog established theories with their EBF $(C, \Psi)$ restrictions
    \item Enable scientific comparison and integration of theories
    \item Document validity domains $\Omega_\mathcal{T}$ for each theory
\end{itemize}

\textbf{In-Scope:}
\begin{itemize}[nosep]
    \item 134 established economic and behavioral theories in 13 categories
    \item EBF restriction mapping $(C_{ij}, \Psi, U^{IDN}, \lambda, \beta)$
    \item Validity domains and known limitations
    \item Key literature references
\end{itemize}

\textbf{Out-of-Scope (delegiert an andere Appendices):}
\begin{itemize}[nosep]
    \item Formal mathematical proofs $\rightarrow$ Appendix UNMAPPED_DER (FORMAL-LIMITS)
    \item Model design workflow $\rightarrow$ Appendix UNMAPPED_DSN (METHOD-DESIGN)
    \item Custom model storage $\rightarrow$ Appendix UNMAPPED_REG (METHOD-REGISTRY)
    \item Parameter values $\rightarrow$ Appendix UNMAPPED_BBB (CORE-WHERE)
\end{itemize}

\textbf{Constraints:}
\begin{itemize}[nosep]
    \item Catalog reflects current scientific consensus
    \item Restrictions are stylized representations
    \item Updates required as new theories emerge
\end{itemize}

\textbf{Lieferobjekte (Deliverables):}
\begin{itemize}[nosep]
    \item 13 Category tables (134 models)
    \item EBF restriction specifications
    \item Validity domain mappings
    \item Cross-reference to formal proofs (DER)
\end{itemize}

\end{tcolorbox}

% =============================================================================
% CROSS-REFERENCE MAP
% =============================================================================
\begin{tcolorbox}[colback=gray!5!white, colframe=gray!75!black,
    title=\textbf{Cross-Reference Map}]

\textbf{Outgoing References:}
\begin{itemize}[nosep]
    \item \textbf{DER (FORMAL-LIMITS):} Formal proofs for 12 key theories
    \item \textbf{DSN (METHOD-DESIGN):} Practical model design workflow
    \item \textbf{BBB (CORE-WHERE):} Parameter values from literature
    \item \textbf{FM (REF-FORMALMAP):} Chapter $\rightarrow$ Proof navigation
    \item \textbf{DP (REF-DEVPIPE):} 6-stage computational pipeline
    \item \textbf{LIT-* Appendices:} Author-specific research integration
\end{itemize}

\textbf{Incoming References:}
\begin{itemize}[nosep]
    \item Chapter 1: Framework introduction
    \item Chapter 12: Theory integration
    \item \textbf{DP Stage 2:} Theory Matching step in pipeline
    \item All scientific users seeking theory comparison
\end{itemize}

\end{tcolorbox}

% =============================================================================
% CONNECTION ARCHITECTURE (Detailed)
% =============================================================================
\begin{tcolorbox}[colback=blue!3!white, colframe=blue!60!black,
    title=\textbf{Connection Architecture: MS $\leftrightarrow$ Other Appendices}]

\textbf{Visual Overview:}

\begin{center}
\begin{tabular}{|c|c|c|c|c|}
\hline
\multicolumn{5}{|c|}{\textbf{Appendix Connection Map}} \\
\hline
& & \textbf{DER} & & \\
& & (Formal Proofs) & & \\
& & $\uparrow$ & & \\
\hline
\textbf{DSN} & $\longleftarrow$ & \textbf{MS} & $\longrightarrow$ & \textbf{BBB} \\
(Design) & & \textit{(134 Theories)} & & (Parameters) \\
\hline
& & $\downarrow$ & & \\
& & \textbf{REG} & & \\
& & (Registry) & & \\
\hline
\multicolumn{2}{|c|}{\textbf{G} (Glossary)} & $\leftrightarrow$ & \multicolumn{2}{c|}{\textbf{LIT-*} (Papers)} \\
\hline
\end{tabular}
\end{center}

\vspace{0.5em}

\textbf{Connection Matrix:}

\begin{center}
\small
\begin{tabular}{|l|c|l|l|}
\hline
\textbf{Appendix} & \textbf{Direction} & \textbf{Purpose} & \textbf{Data Flow} \\
\hline
\textbf{DER} & MS $\rightarrow$ DER & Formal proofs & Theory $\mapsto$ Proof \\
(FORMAL-LIMITS) & DER $\rightarrow$ MS & Validity confirmation & Proof $\mapsto$ $\Omega_\mathcal{T}$ \\
\hline
\textbf{BBB} & MS $\rightarrow$ BBB & Parameter lookup & Theory $\mapsto$ $\Theta$ values \\
(CORE-WHERE) & BBB $\rightarrow$ MS & Empirical bounds & $\lambda, \beta$ ranges \\
\hline
\textbf{DSN} & DSN $\rightarrow$ MS & Theory matching & 10C $\mapsto$ similar theory \\
(METHOD-DESIGN) & MS $\rightarrow$ DSN & Restriction templates & $R_\mathcal{T}$ constraints \\
\hline
\textbf{REG} & REG $\rightarrow$ MS & Validation & Custom $\mapsto$ established \\
(METHOD-REGISTRY) & MS $\rightarrow$ REG & Seed models & Theory $\mapsto$ template \\
\hline
\textbf{G} & MS $\rightarrow$ G & Symbol definitions & $C_{ij}, \Psi, \lambda$ \\
(Glossary) & G $\rightarrow$ MS & Consistent notation & Standard symbols \\
\hline
\textbf{LIT-*} & MS $\rightarrow$ LIT & Paper references & Theory $\mapsto$ \texttt{bib\_keys} \\
(Literature) & LIT $\rightarrow$ MS & New theories & Paper $\mapsto$ \texttt{theory\_support} \\
\hline
\end{tabular}
\end{center}

\end{tcolorbox}

% =============================================================================
% DATA FLOW: Practitioner Workflow
% =============================================================================
\begin{tcolorbox}[colback=green!3!white, colframe=green!60!black,
    title=\textbf{Data Flow: Hybrid Practitioner-Scientist Workflow}]

\begin{center}
\begin{tabular}{ccccccc}
\textbf{Step 1} & & \textbf{Step 2} & & \textbf{Step 3} & & \textbf{Step 4} \\
\hline
10C Design & $\xrightarrow{\text{match}}$ & \textbf{MS} & $\xrightarrow{\text{proof}}$ & DER & $\xrightarrow{\text{params}}$ & BBB \\
(DSN) & & Theory Match & & Formal Proof & & Parameters \\
\end{tabular}
\end{center}

\vspace{0.5em}

\textbf{Workflow Example:}
\begin{enumerate}[nosep]
    \item \textbf{DSN:} Practitioner designs model with $\beta < 1$ (present bias)
    \item \textbf{MS:} Match to Hyperbolic Discounting (MS-TD-002)
    \item \textbf{DER:} Verify formal proof for validity domain
    \item \textbf{BBB:} Get literature-validated $\beta$ range [0.7, 0.95]
    \item \textbf{REG:} Store customized model with theory link
\end{enumerate}

\end{tcolorbox}

% =============================================================================
% BIDIRECTIONAL PAPER INTEGRATION
% =============================================================================
\begin{tcolorbox}[colback=yellow!5!white, colframe=yellow!60!black,
    title=\textbf{Bidirectional Paper Integration}]

\textbf{Connection to Bibliography (\texttt{bcm\_master.bib}):}

Each theory in this catalog has a \texttt{bib\_keys} field linking to papers in the
master bibliography. Conversely, papers can specify \texttt{theory\_support} to
indicate which theories they provide evidence for.

\begin{center}
\begin{tabular}{|l|l|}
\hline
\textbf{Theory Catalog} & \textbf{Bibliography} \\
\hline
\texttt{id: MS-RD-001} & \texttt{@article\{PAP-kahneman1979prospectprospect,} \\
\texttt{name: Prospect Theory} & \texttt{  theory\_support = \{MS-RD-001, MS-RD-002\},} \\
\texttt{bib\_keys: [PAP-kahneman1979prospectprospect]} & \texttt{  ...\}} \\
\hline
\end{tabular}
\end{center}

\textbf{Lookup Scripts:}
\begin{itemize}[nosep]
    \item \texttt{python scripts/theory\_papers.py --theory MS-RD-001} \\
          $\rightarrow$ Find all papers supporting Prospect Theory
    \item \texttt{python scripts/theory\_papers.py --paper PAP-kahneman1979prospectprospect} \\
          $\rightarrow$ Find all theories this paper supports
    \item \texttt{python scripts/theory\_papers.py --match-10c "lambda > 1"} \\
          $\rightarrow$ Match 10C restriction to theories
\end{itemize}

\textbf{Auto-Sync (Pre-Commit Hook):}
When papers with \texttt{theory\_support} are committed, links are automatically
synchronized to this catalog via \texttt{scripts/theory\_learning.py}.

\end{tcolorbox}

% =============================================================================
% CHAPTER LINKAGE
% =============================================================================
\begin{tcolorbox}[colback=purple!5!white, colframe=purple!75!black,
    title=\textbf{Chapter Linkage}]

\textbf{Primary Chapter:} Chapter 12 (Integration of Established Theories)

\textbf{Secondary Chapters:}
\begin{itemize}[nosep]
    \item Chapter 1: Central question (``organizes rather than replaces'')
    \item Chapter 5: Complementarity foundations
\end{itemize}

\end{tcolorbox}

% =============================================================================
% FUNDAMENTAL QUESTION
% =============================================================================
\section{The Fundamental Question}
\label{sec:ms-fundamental}

\begin{tcolorbox}[colback=red!5!white, colframe=red!75!black,
    title=\textbf{The Fundamental Question}]

\textbf{Scientific Question:} How do established economic and behavioral
theories relate to each other within the $(C, \Psi)$ framework?

\textbf{Answer:} Each theory $\mathcal{T}$ corresponds to specific restrictions
$R_\mathcal{T}$ on the general framework. EBF does not refute theories—it
specifies their \textit{domain of validity}:
\begin{equation}
\mathcal{T} \text{ is valid when } \Psi \in \Omega_\mathcal{T}
\end{equation}

\end{tcolorbox}

% =============================================================================
% THE THREE PERSPECTIVES (FUNDAMENTAL LAW)
% =============================================================================
\section{EBF Fundamental Law: Three Perspectives}
\label{sec:ms-fundamental-law}

\begin{tcolorbox}[colback=green!5!white, colframe=green!75!black,
    title=\textbf{Axiom UNMAPPED_MS-0: The Organizing Metatheory Principle}]

EBF is an \textbf{organizing framework}, not a competing theory. It can be
used in three ways:

\begin{enumerate}
    \item \textbf{Practitioner Perspective:} Use 10C dimensions directly to
          design evidence-based models for prediction.
          \begin{itemize}[nosep]
              \item Workflow: Appendix UNMAPPED_DSN (METHOD-DESIGN)
              \item Storage: Appendix UNMAPPED_REG (METHOD-REGISTRY)
          \end{itemize}

    \item \textbf{Scientist Perspective:} Understand how established theories
          relate as special cases in the $(C, \Psi)$ space.
          \begin{itemize}[nosep]
              \item Catalog: \textbf{This Appendix (MS)}
              \item Proofs: Appendix UNMAPPED_DER (FORMAL-LIMITS)
          \end{itemize}

    \item \textbf{Hybrid Perspective:} Design models with 10C, then map to
          established theories for empirical grounding and literature support.
          \begin{itemize}[nosep]
              \item Design: Appendix UNMAPPED_DSN
              \item Grounding: This Appendix (MS)
          \end{itemize}
\end{enumerate}

\textbf{All three are valid.} The choice depends on the user's goals.

\end{tcolorbox}

% =============================================================================
% EBF RESTRICTION SCHEMA
% =============================================================================
\section{EBF Restriction Schema}
\label{sec:ms-schema}

Every established theory $\mathcal{T}$ can be characterized by restrictions on
four dimensions:

\begin{tcolorbox}[colback=blue!5!white, colframe=blue!75!black,
    title=\textbf{The Four Restriction Dimensions}]

\begin{tabular}{@{}lll@{}}
\textbf{Dimension} & \textbf{Symbol} & \textbf{Question} \\
\midrule
Complementarity & $C_{ij}$ & Are utilities interdependent? \\
Context & $\Psi$ & Is context exogenous, endogenous, or fixed? \\
Identity & $U^{IDN}$ & Does identity utility matter? \\
Loss Aversion & $\lambda$ & Is loss aversion present? ($\lambda > 1$) \\
\end{tabular}

\medskip
\textbf{General Framework:}
\begin{align}
U_i &= U_i(x_i, x_{-i}, \Psi) \\
C_{ij} &= \frac{\partial^2 U_i}{\partial x_i \partial x_j} \\
\Psi_{t+1} &= f(\Psi_t, a_t)
\end{align}

\textbf{Each theory imposes restrictions} that simplify this structure.

\end{tcolorbox}

% =============================================================================
% CATEGORY 1: CLASSICAL ECONOMICS
% =============================================================================
\section{Category 1: Classical Economics (12 Models)}
\label{sec:ms-classical}

\begin{tcolorbox}[colback=gray!5!white, colframe=gray!75!black,
    title=\textbf{Classical Economics: The Separability Restriction}]

\textbf{Defining Restriction:} $C_{ij} = 0$ for $i \neq j$ (separable utilities)

\textbf{Validity Domain:} Large anonymous markets, price-taking agents,
no social interaction effects.

\end{tcolorbox}

\small
\begin{longtable}{@{}p{3cm}p{2.5cm}p{4cm}p{4cm}@{}}
\toprule
\textbf{Model} & \textbf{Authors} & \textbf{EBF Restrictions} & \textbf{Validity Domain} \\
\midrule
\endhead

Arrow-Debreu GE & Arrow, Debreu (1954) &
$C_{ij}=0$, $\Psi$ fixed, $U^{IDN}=0$ &
Complete markets, price-taking \\

Expected Utility & von Neumann, Morgenstern (1944) &
$C_{ij}=0$, $\lambda=1$, $\Psi$ fixed &
Consistent preferences, no framing \\

Rational Expectations & Muth (1961), Lucas (1972) &
$C_{ij}=0$, $\Psi$ known, no bias &
Information efficient markets \\

Permanent Income & Friedman (1957) &
$C_{ij}=0$, $\Psi_{income}$ stochastic &
Consumption smoothing, no liquidity constraints \\

Life-Cycle Model & Modigliani (1966) &
$C_{ij}=0$, $\Psi_{age}$ deterministic &
Perfect capital markets, certain lifespan \\

CAPM & Sharpe (1964), Lintner (1965) &
$C_{ij}=\text{Cov}(r_i,r_j)$, $\Psi$ fixed &
Mean-variance preferences, no friction \\

Efficient Markets & Fama (1970) &
$C_{ij}=0$, $\Psi_{info}$ instantaneous &
No arbitrage, rational traders \\

Heckscher-Ohlin & Heckscher (1919), Ohlin (1933) &
$C_{ij}=0$, $\Psi_{factor}$ only varies &
Identical tech \& preferences across countries \\

Ricardian Trade & Ricardo (1817) &
$C_{ij}=0$, $\Psi_{tech}$ varies &
Comparative advantage, no transport costs \\

Solow Growth & Solow (1956) &
$C_{ij}=0$, $\Psi_{tech}$ exogenous &
Diminishing returns, exogenous tech \\

Ramsey Model & Ramsey (1928) &
$C_{ij}=0$, $\Psi$ fixed, optimal saving &
Infinite horizon, perfect foresight \\

Coase Theorem & Coase (1960) &
$C_{ij}=0$, $\Psi_{transaction}=0$ &
Zero transaction costs, clear property rights \\

\bottomrule
\end{longtable}
\normalsize

% =============================================================================
% CATEGORY 2: BEHAVIORAL ECONOMICS - SOCIAL PREFERENCES
% =============================================================================
\section{Category 2: Social Preferences (15 Models)}
\label{sec:ms-social}

\begin{tcolorbox}[colback=gray!5!white, colframe=gray!75!black,
    title=\textbf{Social Preferences: The Interdependence Extension}]

\textbf{Defining Restriction:} $C_{ij} \neq 0$ (utilities are interdependent)

\textbf{Validity Domain:} Small groups, salient interactions,
observable outcomes.

\end{tcolorbox}

\small
\begin{longtable}{@{}p{3cm}p{2.5cm}p{4cm}p{4cm}@{}}
\toprule
\textbf{Model} & \textbf{Authors} & \textbf{EBF Restrictions} & \textbf{Validity Domain} \\
\midrule
\endhead

Fehr-Schmidt & Fehr, Schmidt (1999) &
$C_{ij}=f(\alpha,\beta) \neq 0$, outcome-based &
Lab experiments, salient payoffs \\

Bolton-Ockenfels ERC & Bolton, Ockenfels (2000) &
$C_{ij} \neq 0$, share-based &
Relative standing matters \\

Charness-Rabin & Charness, Rabin (2002) &
$C_{ij} \neq 0$, intention-dependent &
Intentions observable \\

Falk-Fischbacher & Falk, Fischbacher (2006) &
$C_{ij}=f(\text{intentions})$ &
Reciprocity, kindness perception \\

Levine Reciprocity & Levine (1998) &
$C_{ij}=f(\text{type beliefs})$ &
Type-based reciprocity \\

Rabin Fairness & Rabin (1993) &
$C_{ij}=f(\text{kindness})$ &
Psychological game theory \\

Inequality Aversion & Fehr, Schmidt (1999) &
$C_{ij} \neq 0$, $\alpha \geq \beta$ &
Envy stronger than guilt \\

Pure Altruism & Andreoni (1989) &
$C_{ij} > 0$ (warm glow) &
Giving behavior \\

Impure Altruism & Andreoni (1990) &
$C_{ij} > 0$, warm glow + impact &
Charitable giving \\

Strong Reciprocity & Gintis et al. (2003) &
$C_{ij} = f(\text{norm compliance})$ &
Punishment of defectors \\

Conditional Cooperation & Fischbacher et al. (2001) &
$C_{ij} = f(\text{others' contributions})$ &
Public goods games \\

Social Norms & Bicchieri (2006) &
$C_{ij} = f(\Psi_{norm})$ &
Norm salience, expectations \\

Gift Exchange & Akerlof (1982) &
$C_{ij} > 0$ (reciprocity) &
Labor markets, effort-wage \\

Efficiency Wages & Shapiro, Stiglitz (1984) &
$C_{ij} > 0$ (monitoring) &
Principal-agent with shirking \\

Trust Game & Berg et al. (1995) &
$C_{ij} = f(\text{trust}, \text{reciprocity})$ &
Sequential exchange \\

\bottomrule
\end{longtable}
\normalsize

% =============================================================================
% CATEGORY 3: BEHAVIORAL ECONOMICS - PROSPECT THEORY & REFERENCE DEPENDENCE
% =============================================================================
\section{Category 3: Reference Dependence (10 Models)}
\label{sec:ms-reference}

\begin{tcolorbox}[colback=gray!5!white, colframe=gray!75!black,
    title=\textbf{Reference Dependence: The Loss Aversion Extension}]

\textbf{Defining Restriction:} $\lambda > 1$ (losses loom larger than gains)

\textbf{Validity Domain:} Risky choices, status quo effects,
framing manipulations.

\end{tcolorbox}

\small
\begin{longtable}{@{}p{3cm}p{2.5cm}p{4cm}p{4cm}@{}}
\toprule
\textbf{Model} & \textbf{Authors} & \textbf{EBF Restrictions} & \textbf{Validity Domain} \\
\midrule
\endhead

Prospect Theory & Kahneman, Tversky (1979) &
$\lambda \approx 2.25$, $\Psi_{ref}$ endogenous &
Risky choice, framing \\

Cumulative PT & Tversky, Kahneman (1992) &
$\lambda > 1$, probability weighting &
Multi-outcome lotteries \\

Loss Aversion & Kahneman et al. (1991) &
$\lambda > 1$, WTA $>$ WTP &
Endowment effect \\

Reference-Dep. Preferences & Kőszegi, Rabin (2006) &
$\lambda > 1$, $\Psi_{ref}$ = expectations &
Expectations-based reference \\

Endowment Effect & Thaler (1980) &
$\lambda > 1$, ownership shifts $\Psi_{ref}$ &
Ownership, possession \\

Status Quo Bias & Samuelson, Zeckhauser (1988) &
$\lambda > 1$, default = reference &
Default options, inertia \\

Disposition Effect & Shefrin, Statman (1985) &
$\lambda > 1$, purchase price = ref &
Stock selling behavior \\

Mental Accounting & Thaler (1985) &
$\lambda > 1$, separate accounts &
Fungibility violations \\

Narrow Bracketing & Read et al. (1999) &
$\lambda > 1$, isolated decisions &
Sequence vs. portfolio \\

House Money Effect & Thaler, Johnson (1990) &
$\lambda$ varies with prior gains &
Prior outcomes affect risk \\

\bottomrule
\end{longtable}
\normalsize

% =============================================================================
% CATEGORY 4: TIME PREFERENCES
% =============================================================================
\section{Category 4: Time Preferences (10 Models)}
\label{sec:ms-time}

\begin{tcolorbox}[colback=gray!5!white, colframe=gray!75!black,
    title=\textbf{Time Preferences: The Discounting Extension}]

\textbf{Defining Restriction:} Non-exponential discounting, present bias ($\beta < 1$)

\textbf{Validity Domain:} Intertemporal choices, self-control problems.

\end{tcolorbox}

\small
\begin{longtable}{@{}p{3cm}p{2.5cm}p{4cm}p{4cm}@{}}
\toprule
\textbf{Model} & \textbf{Authors} & \textbf{EBF Restrictions} & \textbf{Validity Domain} \\
\midrule
\endhead

Exponential Discounting & Samuelson (1937) &
$\delta^t$, time-consistent &
Long-horizon, no self-control issues \\

Hyperbolic Discounting & Ainslie (1975), Laibson (1997) &
$\beta\delta^t$, $\beta < 1$ &
Present bias, preference reversal \\

Quasi-Hyperbolic ($\beta$-$\delta$) & Laibson (1997) &
$\beta < 1$, two-parameter &
Tractable present bias \\

Naive Present Bias & O'Donoghue, Rabin (1999) &
$\beta < 1$, $\hat{\beta} = 1$ &
Unaware of future self-control \\

Sophisticated Present Bias & O'Donoghue, Rabin (1999) &
$\beta < 1$, $\hat{\beta} = \beta$ &
Aware, seeks commitment \\

Dual-Self Model & Fudenberg, Levine (2006) &
Short-run vs. long-run self &
Self-control as game \\

Temptation \& Self-Control & Gul, Pesendorfer (2001) &
Menu-dependent utility &
Temptation costs \\

Projection Bias & Loewenstein et al. (2003) &
$\Psi_{future}$ underestimated &
Hot-cold empathy gap \\

Habit Formation & Becker, Murphy (1988) &
$\Psi_t = f(\Psi_{t-1}, a_{t-1})$ &
State-dependent preferences \\

Addiction Model & Becker, Murphy (1988) &
$C_{t,t+1} > 0$ (adjacent comp.) &
Rational addiction \\

\bottomrule
\end{longtable}
\normalsize

% =============================================================================
% CATEGORY 5: IDENTITY & BELIEFS
% =============================================================================
\section{Category 5: Identity \& Beliefs (12 Models)}
\label{sec:ms-identity}

\begin{tcolorbox}[colback=gray!5!white, colframe=gray!75!black,
    title=\textbf{Identity \& Beliefs: The $U^{IDN}$ Extension}]

\textbf{Defining Restriction:} $U^{IDN} \neq 0$ (identity utility matters)

\textbf{Validity Domain:} Group membership, self-image, motivated reasoning.

\end{tcolorbox}

\small
\begin{longtable}{@{}p{3cm}p{2.5cm}p{4cm}p{4cm}@{}}
\toprule
\textbf{Model} & \textbf{Authors} & \textbf{EBF Restrictions} & \textbf{Validity Domain} \\
\midrule
\endhead

Identity Economics & Akerlof, Kranton (2000) &
$U^{IDN} \neq 0$, category-based &
Group membership, norms \\

Moral Identity & Bénabou, Tirole (2011) &
$U^{IDN} = f(\mathbb{E}[\theta|a])$ &
Self-signaling, image \\

Self-Signaling & Bodner, Prelec (2003) &
$U^{IDN} = f(\text{diagnostic actions})$ &
Actions reveal character \\

Social Image & Bénabou, Tirole (2006) &
$U^{IDN} = f(\text{public perception})$ &
Reputation, signaling \\

Cognitive Dissonance & Festinger (1957) &
$U^{IDN}$ penalizes inconsistency &
Belief-action alignment \\

Motivated Beliefs & Bénabou, Tirole (2016) &
$U^{IDN}$ shapes belief formation &
Wishful thinking \\

Confirmation Bias & Rabin, Schrag (1999) &
$\Psi_{beliefs}$ sticky to priors &
Asymmetric updating \\

Overconfidence & Moore, Healy (2008) &
$\Psi_{ability}$ biased upward &
Self-assessment errors \\

Bayesian Updating & Bayes (1763) &
$\Psi_{beliefs}$ = posterior &
Rational belief updating \\

Non-Bayesian Updating & Grether (1980) &
$\Psi_{beliefs} \neq$ posterior &
Base rate neglect \\

Cultural Transmission & Bisin, Verdier (2001) &
$U^{IDN}$ transmitted to children &
Intergenerational preferences \\

Group Identity & Tajfel (1970) &
$U^{IDN} = f(\text{in-group status})$ &
Minimal group paradigm \\

\bottomrule
\end{longtable}
\normalsize

% =============================================================================
% CATEGORY 6: INSTITUTIONAL & ORGANIZATIONAL
% =============================================================================
\section{Category 6: Institutions \& Organizations (15 Models)}
\label{sec:ms-institutions}

\begin{tcolorbox}[colback=gray!5!white, colframe=gray!75!black,
    title=\textbf{Institutions: The Endogenous Context Extension}]

\textbf{Defining Restriction:} $\Psi_{t+1} = f(\Psi_t, a_t)$ (context evolves)

\textbf{Validity Domain:} Long-run dynamics, multiple equilibria,
path dependence.

\end{tcolorbox}

\small
\begin{longtable}{@{}p{3cm}p{2.5cm}p{4cm}p{4cm}@{}}
\toprule
\textbf{Model} & \textbf{Authors} & \textbf{EBF Restrictions} & \textbf{Validity Domain} \\
\midrule
\endhead

Acemoglu-Robinson & Acemoglu, Robinson (2012) &
$\Psi_{inst}$ endogenous, multiple eq. &
Long-run development \\

North Institutions & North (1990) &
$\Psi_{inst}$ = formal + informal rules &
Transaction costs, history \\

Williamson TCE & Williamson (1985) &
$\Psi_{transaction}$ determines boundaries &
Make vs. buy \\

Roberts Org Design & Roberts (2004) &
$C_{ij} > 0$ (supermodular) &
Organizational fit \\

Milgrom-Roberts & Milgrom, Roberts (1990) &
$C_{ij} > 0$, strategic comp. &
Modern manufacturing \\

Aghion-Howitt & Aghion, Howitt (1992) &
$\Psi_{tech}$ endogenous via R\&D &
Creative destruction \\

Romer Growth & Romer (1990) &
$\Psi_{ideas}$ endogenous, non-rival &
Idea-based growth \\

Hall-Soskice VoC & Hall, Soskice (2001) &
$\Psi_{inst}$ = LME vs. CME &
Comparative advantage \\

Ostrom Commons & Ostrom (1990) &
$\Psi_{rules}$ community-governed &
Common pool resources \\

Greif Institutions & Greif (2006) &
$\Psi_{beliefs}$ self-enforcing &
Historical institutions \\

Jensen-Meckling & Jensen, Meckling (1976) &
$C_{ij} = f(\text{agency costs})$ &
Principal-agent \\

Holmström Contracts & Holmström (1979) &
$C_{ij} = f(\text{signal precision})$ &
Moral hazard \\

Grossman-Hart & Grossman, Hart (1986) &
$\Psi_{ownership}$ = residual control &
Property rights \\

Hart-Moore & Hart, Moore (1990) &
$\Psi_{contract}$ incomplete &
Incomplete contracts \\

Tirole Regulation & Tirole (1988) &
$\Psi_{info}$ asymmetric &
Regulatory design \\

\bottomrule
\end{longtable}
\normalsize

% =============================================================================
% CATEGORY 7: INFORMATION & LEARNING
% =============================================================================
\section{Category 7: Information \& Learning (12 Models)}
\label{sec:ms-information}

\begin{tcolorbox}[colback=gray!5!white, colframe=gray!75!black,
    title=\textbf{Information: The Asymmetry Extension}]

\textbf{Defining Restriction:} $\Psi_{info}$ asymmetric or costly

\textbf{Validity Domain:} Markets with hidden information, signaling, screening.

\end{tcolorbox}

\small
\begin{longtable}{@{}p{3cm}p{2.5cm}p{4cm}p{4cm}@{}}
\toprule
\textbf{Model} & \textbf{Authors} & \textbf{EBF Restrictions} & \textbf{Validity Domain} \\
\midrule
\endhead

Akerlof Lemons & Akerlof (1970) &
$\Psi_{quality}$ private info &
Adverse selection \\

Spence Signaling & Spence (1973) &
$\Psi_{type}$ signaled via costly action &
Education as signal \\

Stiglitz Screening & Rothschild, Stiglitz (1976) &
$\Psi_{type}$ revealed via menu &
Insurance markets \\

Herd Behavior & Banerjee (1992) &
$\Psi_{beliefs}$ = f(others' actions) &
Information cascades \\

Global Games & Morris, Shin (2003) &
$\Psi_{signal}$ noisy, coordination &
Speculative attacks \\

Cheap Talk & Crawford, Sobel (1982) &
$\Psi_{message}$ costless, bias &
Strategic communication \\

Persuasion & Kamenica, Gentzkow (2011) &
$\Psi_{info}$ designed by sender &
Bayesian persuasion \\

Rational Inattention & Sims (2003) &
$\Psi_{attention}$ limited capacity &
Information processing costs \\

Salience Theory & Bordalo et al. (2012) &
$\Psi_{salient}$ attracts attention &
Context-dependent attention \\

Gabaix Sparsity & Gabaix (2014) &
$\Psi_{complex}$ simplified &
Sparse representations \\

Learning Models & Various &
$\Psi_{beliefs,t+1} = f(\Psi_t, \text{data})$ &
Belief dynamics \\

Experience Effects & Malmendier, Nagel (2011) &
$\Psi_{beliefs}$ = f(personal history) &
Depression babies \\

\bottomrule
\end{longtable}
\normalsize

% =============================================================================
% CATEGORY 8: GAME THEORY & STRATEGIC INTERACTION
% =============================================================================
\section{Category 8: Strategic Interaction (14 Models)}
\label{sec:ms-games}

\begin{tcolorbox}[colback=gray!5!white, colframe=gray!75!black,
    title=\textbf{Strategic Interaction: The Multiple Equilibria Extension}]

\textbf{Defining Restriction:} $C_{ij} \neq 0$, strategic interdependence

\textbf{Validity Domain:} Small numbers, strategic awareness, repeated play.

\end{tcolorbox}

\small
\begin{longtable}{@{}p{3cm}p{2.5cm}p{4cm}p{4cm}@{}}
\toprule
\textbf{Model} & \textbf{Authors} & \textbf{EBF Restrictions} & \textbf{Validity Domain} \\
\midrule
\endhead

Nash Equilibrium & Nash (1950) &
$C_{ij} \neq 0$, best responses &
Rational players, common knowledge \\

Subgame Perfect & Selten (1965) &
$C_{ij} \neq 0$, backward induction &
Sequential games \\

Bayesian Nash & Harsanyi (1967) &
$C_{ij} \neq 0$, $\Psi_{type}$ private &
Incomplete information \\

Quantal Response & McKelvey, Palfrey (1995) &
$C_{ij} \neq 0$, noisy best response &
Bounded rationality in games \\

Level-k Thinking & Nagel (1995) &
$C_{ij} \neq 0$, limited depth &
Beauty contests \\

Cognitive Hierarchy & Camerer et al. (2004) &
$C_{ij} \neq 0$, Poisson levels &
Strategic sophistication \\

Focal Points & Schelling (1960) &
$C_{ij} \neq 0$, $\Psi_{salient}$ coordinates &
Coordination games \\

Evolutionary Game Theory & Maynard Smith (1982) &
$C_{ij} \neq 0$, population dynamics &
Long-run stability \\

Repeated Games & Fudenberg, Maskin (1986) &
$C_{ij} \neq 0$, folk theorem &
Sustained cooperation \\

Mechanism Design & Myerson (1981) &
$C_{ij}$ designed by principal &
Optimal auctions \\

Matching Markets & Gale, Shapley (1962) &
$C_{ij} \neq 0$, stability &
Two-sided matching \\

Market Design & Roth (2002) &
$\Psi_{rules}$ designed &
Kidney exchange, schools \\

Bargaining & Nash (1950), Rubinstein (1982) &
$C_{ij} \neq 0$, surplus division &
Bilateral negotiation \\

Auction Theory & Vickrey (1961) &
$C_{ij} \neq 0$, $\Psi_{value}$ private &
Revenue maximization \\

\bottomrule
\end{longtable}
\normalsize

% =============================================================================
% CATEGORY 9: BEHAVIORAL FINANCE
% =============================================================================
\section{Category 9: Behavioral Finance (8 Models)}
\label{sec:ms-behfin}

\begin{tcolorbox}[colback=gray!5!white, colframe=gray!75!black,
    title=\textbf{Behavioral Finance: Market Anomalies Extension}]

\textbf{Defining Restriction:} $\lambda > 1$ + limits to arbitrage + noise traders

\textbf{Validity Domain:} Financial markets, asset pricing, investor behavior.

\end{tcolorbox}

\small
\begin{longtable}{@{}p{3cm}p{2.5cm}p{4cm}p{4cm}@{}}
\toprule
\textbf{Model} & \textbf{Authors} & \textbf{EBF Restrictions} & \textbf{Validity Domain} \\
\midrule
\endhead

Limits to Arbitrage & Shleifer, Vishny (1997) &
$\Psi_{risk}$ constrains correction &
Mispricing persists \\

Noise Trader Risk & DeLong et al. (1990) &
$C_{ij} \neq 0$, noise affects prices &
Irrational traders matter \\

Behavioral Portfolio & Shefrin, Statman (2000) &
$\lambda > 1$, mental layers &
Layered portfolio construction \\

Equity Premium Puzzle & Mehra, Prescott (1985) &
$\lambda \gg 1$ explains premium &
Historical excess returns \\

Dividend Puzzle & Shefrin, Statman (1984) &
$\lambda > 1$, mental accounting &
Preference for dividends \\

Momentum \& Reversal & Jegadeesh, Titman (1993) &
$\Psi_{beliefs}$ slow updating &
Short-term momentum \\

Overreaction & DeBondt, Thaler (1985) &
$\Psi_{beliefs}$ overshoot &
Long-term reversal \\

Home Bias & French, Poterba (1991) &
$\Psi_{familiar}$ preference &
Domestic overweighting \\

\bottomrule
\end{longtable}
\normalsize

% =============================================================================
% CATEGORY 10: NUDGE & CHOICE ARCHITECTURE
% =============================================================================
\section{Category 10: Nudge \& Choice Architecture (8 Models)}
\label{sec:ms-nudge}

\begin{tcolorbox}[colback=gray!5!white, colframe=gray!75!black,
    title=\textbf{Nudge: The Context Design Extension}]

\textbf{Defining Restriction:} $\Psi$ is designable; behavior responds to architecture

\textbf{Validity Domain:} Policy design, defaults, framing interventions.

\end{tcolorbox}

\small
\begin{longtable}{@{}p{3cm}p{2.5cm}p{4cm}p{4cm}@{}}
\toprule
\textbf{Model} & \textbf{Authors} & \textbf{EBF Restrictions} & \textbf{Validity Domain} \\
\midrule
\endhead

Libertarian Paternalism & Thaler, Sunstein (2003) &
$\Psi_{default}$ designable &
Freedom-preserving interventions \\

Default Effects & Madrian, Shea (2001) &
$\Psi_{default}$ = reference point &
Retirement savings, organ donation \\

MINDSPACE Framework & Dolan et al. (2012) &
9 $\Psi$ channels for behavior change &
UK policy applications \\

EAST Framework & BIT (2014) &
Easy, Attractive, Social, Timely &
Simplified nudge design \\

Boost vs. Nudge & Hertwig, Grüne-Yanoff (2017) &
$\Psi_{competence}$ vs. $\Psi_{architecture}$ &
Competence building vs. steering \\

Sludge & Sunstein (2020) &
$\Psi_{friction}$ as barrier &
Harmful choice architecture \\

Active Choosing & Carroll et al. (2009) &
$\Psi_{default} = \emptyset$, forced choice &
When defaults inappropriate \\

Smart Defaults & Goldstein et al. (2008) &
$\Psi_{default}$ personalized &
Adaptive defaults \\

\bottomrule
\end{longtable}
\normalsize

% =============================================================================
% CATEGORY 11: NEUROECONOMICS & DUAL SYSTEMS
% =============================================================================
\section{Category 11: Neuroeconomics \& Dual Systems (6 Models)}
\label{sec:ms-neuro}

\begin{tcolorbox}[colback=gray!5!white, colframe=gray!75!black,
    title=\textbf{Neuroeconomics: The Neural Implementation Extension}]

\textbf{Defining Restriction:} Dual systems (S1/S2); neural constraints on $\Psi$ processing

\textbf{Validity Domain:} Cognitive load, automaticity, emotional processing.

\end{tcolorbox}

\small
\begin{longtable}{@{}p{3cm}p{2.5cm}p{4cm}p{4cm}@{}}
\toprule
\textbf{Model} & \textbf{Authors} & \textbf{EBF Restrictions} & \textbf{Validity Domain} \\
\midrule
\endhead

Dual-System (S1/S2) & Kahneman (2011) &
$\Psi_{load}$ determines system &
Fast/slow thinking \\

Somatic Marker & Damasio (1994) &
$U^{E}$ via bodily signals &
Emotion in decision \\

Reward Prediction Error & Schultz (1997) &
$\Psi_{expected}$ vs. $\Psi_{actual}$ &
Dopamine, learning \\

Hot-Cold Empathy Gap & Loewenstein (2005) &
$\Psi_{visceral}$ underestimated &
Projection bias \\

Ego Depletion & Baumeister et al. (1998) &
$\Psi_{willpower}$ limited resource &
Self-control failures \\

Construal Level & Trope, Liberman (2010) &
$\Psi_{distance}$ affects abstraction &
Psychological distance \\

\bottomrule
\end{longtable}
\normalsize

% =============================================================================
% CATEGORY 12: WELLBEING & HAPPINESS
% =============================================================================
\section{Category 12: Wellbeing \& Happiness (6 Models)}
\label{sec:ms-wellbeing}

\begin{tcolorbox}[colback=gray!5!white, colframe=gray!75!black,
    title=\textbf{Wellbeing: The Experienced Utility Extension}]

\textbf{Defining Restriction:} $U^{experienced} \neq U^{decision}$; adaptation effects

\textbf{Validity Domain:} Life satisfaction, hedonic adaptation, meaning.

\end{tcolorbox}

\small
\begin{longtable}{@{}p{3cm}p{2.5cm}p{4cm}p{4cm}@{}}
\toprule
\textbf{Model} & \textbf{Authors} & \textbf{EBF Restrictions} & \textbf{Validity Domain} \\
\midrule
\endhead

Easterlin Paradox & Easterlin (1974) &
$\Psi_{reference}$ = social comparison &
Income vs. happiness \\

Hedonic Treadmill & Brickman, Campbell (1971) &
$\Psi_{adaptation}$ returns to baseline &
Adaptation to events \\

Set Point Theory & Lykken, Tellegen (1996) &
$U^{SWB}$ genetically bounded &
Stable happiness range \\

PERMA Model & Seligman (2011) &
$U = f(P, E, R, M, A)$ &
Positive psychology \\

Peak-End Rule & Kahneman et al. (1993) &
$U^{remembered} = f(\text{peak}, \text{end})$ &
Memory vs. experience \\

Focusing Illusion & Schkade, Kahneman (1998) &
$\Psi_{salient}$ overweighted &
Affective forecasting errors \\

\bottomrule
\end{longtable}
\normalsize

% =============================================================================
% CATEGORY 13: DEVELOPMENT & POVERTY
% =============================================================================
\section{Category 13: Development \& Poverty (6 Models)}
\label{sec:ms-development}

\begin{tcolorbox}[colback=gray!5!white, colframe=gray!75!black,
    title=\textbf{Development: The Scarcity Extension}]

\textbf{Defining Restriction:} $\Psi_{scarcity}$ constrains bandwidth; poverty traps

\textbf{Validity Domain:} Low-income contexts, resource constraints, development.

\end{tcolorbox}

\small
\begin{longtable}{@{}p{3cm}p{2.5cm}p{4cm}p{4cm}@{}}
\toprule
\textbf{Model} & \textbf{Authors} & \textbf{EBF Restrictions} & \textbf{Validity Domain} \\
\midrule
\endhead

Scarcity Model & Mullainathan, Shafir (2013) &
$\Psi_{bandwidth}$ reduced by scarcity &
Cognitive tax of poverty \\

Poverty Traps & Banerjee, Duflo (2011) &
$\Psi_{t+1} = f(\Psi_t)$, multiple equilibria &
S-shaped production \\

Aspirations Window & Ray (2006) &
$\Psi_{aspiration}$ affects effort &
Too far = give up \\

Hope \& Aspirations & Lybbert, Wydick (2018) &
$U^{IDN}$ via future self &
Hope as capital \\

Cash Transfers & Haushofer, Shapiro (2016) &
$\Psi_{liquidity}$ relaxed &
Unconditional transfers \\

Graduation Programs & Banerjee et al. (2015) &
$\Psi$ multi-dimensional intervention &
Big push out of poverty \\

\bottomrule
\end{longtable}
\normalsize

% =============================================================================
% RESULTS: KEY INSIGHTS FROM THE CATALOG
% =============================================================================
\section{Results: Key Insights from the Theory Space}
\label{sec:ms-results}

The comprehensive catalog reveals several important patterns:

\subsection{Restriction Hierarchy}

\begin{enumerate}
    \item \textbf{Classical Economics} ($C_{ij} = 0$, $\Psi$ fixed) is the most
          restrictive configuration—valid only in large anonymous markets.
    \item \textbf{Behavioral Extensions} relax one restriction at a time:
          Social preferences relax $C_{ij} = 0$; Prospect Theory adds $\lambda > 1$;
          Identity Economics adds $U^{IDN} \neq 0$.
    \item \textbf{Institutional Economics} relaxes $\Psi$ fixed, enabling
          endogenous context and multiple equilibria.
    \item \textbf{Full EBF} relaxes all restrictions simultaneously.
\end{enumerate}

\subsection{Coverage Statistics}

\begin{itemize}
    \item \textbf{134 models} across \textbf{13 categories}
    \item \textbf{8 categories} have formal proofs in Appendix UNMAPPED_DER
    \item \textbf{5 newer categories} (Behavioral Finance, Nudge, Neuro, Wellbeing,
          Development) await formal treatment
    \item All models specify \textbf{validity domains} $\Omega_\mathcal{T}$
\end{itemize}

% =============================================================================
% WORKED EXAMPLE
% =============================================================================
\section{Worked Example: Matching a 10C Model to Established Theories}
\label{sec:ms-worked-example}

\begin{tcolorbox}[colback=green!5!white, colframe=green!75!black,
    title=\textbf{Worked Example: Retirement Savings Intervention}]

\textbf{Scenario:} A practitioner designs a retirement savings intervention
using the 10C framework (/design-model workflow).

\textbf{Step 1: 10C Model Specification}
\begin{itemize}[nosep]
    \item WHO: Individual employees (L0)
    \item WHAT: Financial utility $U^F$ primary
    \item HOW: $\gamma_{present,future} < 0$ (present bias)
    \item WHEN: $\Psi_{default}$ = current enrollment status
    \item AWARE: Low awareness of compound interest ($A \approx 0.3$)
    \item READY: High willingness but low action ($W_{base} = 0.7$, $\theta = 0.2$)
\end{itemize}

\textbf{Step 2: Catalog Lookup}

The practitioner searches this catalog for matching restrictions:
\begin{itemize}[nosep]
    \item $\gamma_{present,future} < 0$ $\rightarrow$ \textbf{Category 4: Time Preferences}
    \item $\Psi_{default}$ matters $\rightarrow$ \textbf{Category 10: Nudge}
    \item Low awareness $\rightarrow$ \textbf{Category 7: Information}
\end{itemize}

\textbf{Step 3: Theory Matching}

\begin{center}
\begin{tabular}{@{}lll@{}}
\toprule
\textbf{10C Element} & \textbf{Matched Theory} & \textbf{Literature} \\
\midrule
$\beta < 1$ (present bias) & Quasi-Hyperbolic ($\beta$-$\delta$) & Laibson (1997) \\
$\Psi_{default}$ & Default Effects & Madrian \& Shea (2001) \\
Low $A$ & Rational Inattention & Sims (2003) \\
\bottomrule
\end{tabular}
\end{center}

\textbf{Step 4: Parameter Extraction}

From the matched literature:
\begin{itemize}[nosep]
    \item $\beta \approx 0.7$ (Laibson 1997: range 0.6-0.8)
    \item Default effect: +30-50pp enrollment (Madrian \& Shea 2001)
    \item Attention cost: $\kappa \approx 0.1$ bits/decision (Sims 2003)
\end{itemize}

\textbf{Step 5: Intervention Design}

Based on the matched theories:
\begin{enumerate}[nosep]
    \item Use \textbf{automatic enrollment} (Default Effects)
    \item Provide \textbf{simplified information} (Rational Inattention)
    \item Offer \textbf{commitment devices} (Sophisticated Present Bias)
\end{enumerate}

\textbf{Result:} The practitioner now has a 10C model grounded in 3 established
theories with empirically validated parameters.

\end{tcolorbox}

% =============================================================================
% SUMMARY: THEORY SPACE
% =============================================================================
\section{Summary: The Theory Space}
\label{sec:ms-summary}

\begin{tcolorbox}[colback=blue!5!white, colframe=blue!75!black,
    title=\textbf{Theory Space Overview}]

\textbf{Total Models Cataloged:} 134 (13 categories)

\begin{center}
\begin{tabular}{@{}llcc@{}}
\toprule
\textbf{Category} & \textbf{Key Restriction} & \textbf{Models} & \textbf{Formal Proofs} \\
\midrule
1. Classical Economics & $C_{ij} = 0$ & 12 & DER §3.1 \\
2. Social Preferences & $C_{ij} \neq 0$ & 15 & DER §3.2-3.3 \\
3. Reference Dependence & $\lambda > 1$ & 10 & DER §3.14 \\
4. Time Preferences & $\beta < 1$ & 10 & --- \\
5. Identity \& Beliefs & $U^{IDN} \neq 0$ & 12 & DER §3.4-3.5 \\
6. Institutions & $\Psi$ endogenous & 15 & DER §3.6-3.9 \\
7. Information & $\Psi_{info}$ asymmetric & 12 & --- \\
8. Strategic Interaction & Multiple equilibria & 14 & DER §3.10-3.11 \\
9. Behavioral Finance & Limits to arbitrage & 8 & --- \\
10. Nudge \& Choice Arch. & $\Psi$ designable & 8 & --- \\
11. Neuroeconomics & Dual systems (S1/S2) & 6 & --- \\
12. Wellbeing & $U^{exp} \neq U^{dec}$ & 6 & --- \\
13. Development & $\Psi_{scarcity}$ & 6 & --- \\
\midrule
\textbf{Total} & & \textbf{134} & \\
\bottomrule
\end{tabular}
\end{center}

\end{tcolorbox}

% =============================================================================
% HOW TO USE THIS CATALOG
% =============================================================================
\section{How to Use This Catalog}
\label{sec:ms-usage}

\subsection{For Scientists}

\begin{enumerate}
    \item \textbf{Theory Comparison:} Use restriction columns to compare
          what different theories assume.
    \item \textbf{Validity Domains:} Check when each theory applies
          (and when it doesn't).
    \item \textbf{Literature Grounding:} Find key papers for each model.
    \item \textbf{Formal Proofs:} See Appendix UNMAPPED_DER for mathematical derivations.
\end{enumerate}

\subsection{For Hybrid Users (Practitioner + Scientist)}

\begin{enumerate}
    \item \textbf{Design model} using 10C (Appendix UNMAPPED_DSN).
    \item \textbf{Match to catalog:} Find similar models in this appendix.
    \item \textbf{Extract parameters:} Use literature values from matched models.
    \item \textbf{Cite evidence:} Reference established papers for credibility.
\end{enumerate}

% =============================================================================
% GLOSSARY
% =============================================================================
\section{Glossary}
\label{sec:ms-glossary}

\begin{description}
    \item[$C_{ij}$] Cross-partial derivative $\partial^2 U_i / \partial x_i \partial x_j$;
                    measures utility interdependence.
    \item[$\Psi$] Context vector; situational factors affecting behavior.
    \item[$U^{IDN}$] Identity utility; value from self-image and group membership.
    \item[$\lambda$] Loss aversion parameter; $\lambda > 1$ means losses hurt more than gains help.
    \item[$\beta$] Present bias parameter; $\beta < 1$ means immediacy premium.
    \item[$\Omega_\mathcal{T}$] Validity domain; context configurations where theory $\mathcal{T}$ applies.
    \item[$R_\mathcal{T}$] Restriction set; parametric assumptions that define theory $\mathcal{T}$.
\end{description}

See Appendix G (REF-GLOSSARY) for complete symbol definitions.

% =============================================================================
% OPEN ISSUES
% =============================================================================
\section{Open Issues and Research Agenda}
\label{sec:ms-open-issues}

\begin{enumerate}
    \item \textbf{Formal Proofs:} Categories 9-13 (Behavioral Finance, Nudge,
          Neuroeconomics, Wellbeing, Development) lack formal derivations in
          Appendix UNMAPPED_DER. Future work should extend the limiting case proofs.

    \item \textbf{Parameter Ranges:} Many theories have context-dependent
          parameters (e.g., $\lambda$ varies 1.5-2.5 across domains). A systematic
          meta-analysis could provide domain-specific ranges.

    \item \textbf{Theory Interactions:} The catalog treats theories independently,
          but many interact. For example, present bias ($\beta < 1$) combined with
          loss aversion ($\lambda > 1$) creates specific patterns not captured
          by either alone.

    \item \textbf{Emerging Theories:} New models (e.g., AI decision-making,
          algorithmic nudging) will require ongoing catalog updates.

    \item \textbf{Cross-Cultural Validity:} Most theories were developed in
          WEIRD contexts. Validity domains may need cultural modifiers.
\end{enumerate}

% =============================================================================
% REFERENCES
% =============================================================================
\section{References}
\label{sec:ms-references}

This appendix integrates findings from the master bibliography.
See \texttt{bcm\_master.bib} for complete citations.

\nocite{bcm_master}

\label{sec:ms-footer}
